\documentclass[11pt]{article}
\usepackage[margin=1in]{geometry}
\usepackage[T1]{fontenc}
\usepackage{lmodern,microtype,hyperref,enumitem,amsmath,amssymb,xcolor,booktabs}
\usepackage[numbers,sort&compress]{natbib}
\usepackage{url}
\hypersetup{colorlinks=true,urlcolor=blue,citecolor=blue,linkcolor=blue}

\title{Q11 --- Discussion: Why Do Models Hallucinate, and Can They Know When They Are Unsure?}
\author{DLS 2026 --- Graduate AI-Internals Module}
\date{}

\begin{document}
\maketitle

\section*{Framing the Question}

A student asks why large language models (LLMs) produce confident-sounding
falsehoods, and whether a model can represent its own uncertainty. These are
two distinct but entangled problems. The first concerns the \emph{generative}
tendency of models to emit fluent but ungrounded text. The second concerns
\emph{introspective calibration}: whether the model's internal or externalized
confidence tracks the probability that its output is correct.

It is useful to separate several mechanisms often lumped under
``hallucination.'' Some errors arise from pretraining objectives that reward
plausible continuations rather than truthful ones; \citet{ji2023survey}
provide a taxonomy distinguishing intrinsic from extrinsic hallucination.
Others arise from the statistical structure of the training distribution:
\citet{kalai2024calibrated} argue that, under a cross-entropy objective, a
nontrivial hallucination rate on rare facts is a mathematically expected
consequence of being well-calibrated on the language-modeling task, linking
generation errors to the base rate of ``singleton'' facts seen once in
training. Instruction tuning and RLHF can further amplify confident assertion:
\citet{ouyang2022training} show alignment reshapes output distributions in ways
that do not directly optimize factual accuracy, and reward models may
inadvertently prefer assertive answers.

On the question of self-knowledge, the evidence is mixed. \citet{kadavath2022know}
report that models can be trained or prompted to predict whether they ``know''
an answer, with reasonable calibration on some tasks. \citet{guo2017calibration}
established that modern neural networks are frequently \emph{miscalibrated},
typically overconfident, and that post-hoc methods such as temperature scaling
help. More recent work on verbalized confidence \citep{lin2022teaching} and
on semantic-entropy-based uncertainty estimation \citep{kuhn2023semantic} and
\citep{farquhar2024detecting} suggests that uncertainty signals exist but must
be read out carefully, since token-level probabilities conflate linguistic
and epistemic variability. Retrieval augmentation \citep{lewis2020retrieval}
reduces but does not eliminate the problem, since models may still ignore or
misweight retrieved evidence.

\section*{Five Open Questions for Discussion}

\begin{enumerate}[leftmargin=*,label=\textbf{Q\arabic*.}]

  \item \textbf{Is hallucination a bug or a predictable equilibrium of the
  objective?} If \citet{kalai2024calibrated} are right that error on rare facts
  follows from calibrated density estimation, then eliminating hallucination
  entirely may require trading away calibration or abandoning pure
  next-token likelihood. What objective modifications (e.g., abstention-aware
  losses, selective prediction) could decouple fluency from unwarranted
  assertion without degrading general capability? Can you design a toy
  distribution where this trade-off is provably sharp?

  \item \textbf{Does a model ``know'' it is uncertain, or only know how to
  say so?} \citet{kadavath2022know} and \citet{lin2022teaching} show models
  can emit calibrated-looking confidence statements, but this may reflect
  learned surface correlations rather than access to an internal epistemic
  state. What experiment would distinguish genuine internal uncertainty
  representation (e.g., probes on hidden states, semantic entropy over
  sampled generations per \citet{kuhn2023semantic}) from mere stylistic
  hedging learned from training data?

  \item \textbf{Which layer of the stack is responsible?} Hallucination could
  originate in pretraining distribution gaps, in the decoding procedure, or in
  the RLHF reward model's preference for assertiveness \citep{ouyang2022training}.
  How would you attribute a specific false claim to one of these sources? Can
  controlled ablations (fixed base model, varying decoding temperature,
  varying reward model) isolate the dominant contributor for a given task?

  \item \textbf{Is uncertainty best measured in probability space or semantic
  space?} Token-level entropy can be high simply because many paraphrases are
  valid, not because the model is unsure of the fact. \citet{farquhar2024detecting}
  and \citet{kuhn2023semantic} cluster semantically equivalent generations to
  recover epistemic uncertainty. When does semantic entropy outperform
  log-probability-based confidence, and what are its failure modes when the
  semantic clustering model itself is unreliable?

  \item \textbf{Does grounding transfer calibration?} Retrieval augmentation
  \citep{lewis2020retrieval} supplies external evidence, but a model may still
  hallucinate when evidence is absent, contradictory, or ignored. Can we build
  models whose confidence is explicitly conditioned on retrieval support, such
  that abstention rises when the context is uninformative? How would you
  evaluate whether post-hoc calibration methods \citep{guo2017calibration}
  compose with retrieval, or whether grounding and calibration must be
  co-trained?

\end{enumerate}

\section*{Pointers for the Seminar}

Students should come prepared to argue a position on at least two of the five
questions, grounding claims in the taxonomy of \citet{ji2023survey} and the
empirical self-knowledge results of \citet{kadavath2022know}. A productive
exercise is to reproduce a small semantic-entropy detector following
\citet{farquhar2024detecting} and report where it disagrees with raw
log-probability confidence.

\begin{thebibliography}{99}

\bibitem[Ji et al.(2023)]{ji2023survey}
Z.~Ji, N.~Lee, R.~Frieske, T.~Yu, D.~Su, Y.~Xu, E.~Ishii, Y.~J. Bang,
A.~Madotto, and P.~Fung.
\newblock Survey of hallucination in natural language generation.
\newblock \emph{ACM Computing Surveys}, 55(12):1--38, 2023.

\bibitem[Kalai and Vempala(2024)]{kalai2024calibrated}
A.~T. Kalai and S.~S. Vempala.
\newblock Calibrated language models must hallucinate.
\newblock In \emph{Proceedings of the 56th Annual ACM Symposium on Theory of
Computing (STOC)}, 2024.

\bibitem[Ouyang et al.(2022)]{ouyang2022training}
L.~Ouyang, J.~Wu, X.~Jiang, D.~Almeida, C.~L. Wainwright, P.~Mishkin, C.~Zhang,
S.~Agarwal, K.~Slama, A.~Ray, et al.
\newblock Training language models to follow instructions with human feedback.
\newblock In \emph{Advances in Neural Information Processing Systems (NeurIPS)},
2022.

\bibitem[Kadavath et al.(2022)]{kadavath2022know}
S.~Kadavath, T.~Conerly, A.~Askell, T.~Henighan, D.~Drain, E.~Perez,
N.~Schiefer, Z.~Hatfield-Dodds, N.~DasSarma, E.~Tran-Johnson, et al.
\newblock Language models (mostly) know what they know.
\newblock \emph{arXiv preprint arXiv:2207.05221}, 2022.

\bibitem[Guo et al.(2017)]{guo2017calibration}
C.~Guo, G.~Pleiss, Y.~Sun, and K.~Q. Weinberger.
\newblock On calibration of modern neural networks.
\newblock In \emph{Proceedings of the 34th International Conference on Machine
Learning (ICML)}, 2017.

\bibitem[Lin et al.(2022)]{lin2022teaching}
S.~Lin, J.~Hilton, and O.~Evans.
\newblock Teaching models to express their uncertainty in words.
\newblock \emph{Transactions on Machine Learning Research (TMLR)}, 2022.

\bibitem[Kuhn et al.(2023)]{kuhn2023semantic}
L.~Kuhn, Y.~Gal, and S.~Farquhar.
\newblock Semantic uncertainty: Linguistic invariances for uncertainty
estimation in natural language generation.
\newblock In \emph{International Conference on Learning Representations (ICLR)},
2023.

\bibitem[Farquhar et al.(2024)]{farquhar2024detecting}
S.~Farquhar, J.~Kossen, L.~Kuhn, and Y.~Gal.
\newblock Detecting hallucinations in large language models using semantic
entropy.
\newblock \emph{Nature}, 630:625--630, 2024.

\bibitem[Lewis et al.(2020)]{lewis2020retrieval}
P.~Lewis, E.~Perez, A.~Piktus, F.~Petroni, V.~Karpukhin, N.~Goyal,
H.~K\"uttler, M.~Lewis, W.~Yih, T.~Rockt\"aschel, S.~Riedel, and D.~Kiela.
\newblock Retrieval-augmented generation for knowledge-intensive NLP tasks.
\newblock In \emph{Advances in Neural Information Processing Systems (NeurIPS)},
2020.

\end{thebibliography}

\end{document}